% PRISMA 2020 para resúmenes (ítems 1--12): título (main.tex); objetivos; criterios de elegibilidad, fuentes, riesgo de sesgo y métodos de síntesis (Método); estudios incluidos y resultados (Resultados); limitaciones; interpretación (Conclusiones); financiamiento y registro.
% Nota: el identificador DOI del registro es un identificador, no una cifra del análisis; no existe macro para él.
\begin{abstract}
\textit{Contexto:} En el \en{Text-to-SQL} conversacional cada pregunta depende de los turnos anteriores, y no se identificó una revisión sistemática que analice de forma conjunta la restauración del contexto, el estado del diálogo, la vinculación de esquemas y la autocorrección.
\textit{Objetivo:} Clasificar las técnicas multiturno, contrastarlas con las de un turno y sintetizar su efecto, su evaluación y las brechas.
\textit{Método:} Revisión sistemática según PRISMA 2020 con protocolo registrado. Se consultaron \nBases{} fuentes (Scopus, IEEE Xplore, ACM Digital Library, ACL Anthology y arXiv) el \fechaBusqueda{}, más \en{snowballing}. Se incluyeron estudios empíricos multiturno y, de un turno, solo los que tenían como contribución central la vinculación de esquemas o la autocorrección. Modelos de IA cribaron y extrajeron datos con verificación de la autora; la calidad se evaluó con \nItemsCalidad{} ítems al estilo de Kitchenham; la síntesis fue narrativa (SWiM).
\textit{Resultados:} Se incluyeron \nEstudios{} estudios, \nMultiturno{} de ellos multiturno. La proporción multiturno (D4), descriptiva y dentro del marco de elegibilidad, fue \DcuatroPct~\% (IC 95~\% de \DcuatroICinf{} a \DcuatroICsup~\%); la regla registrada la clasifica como «\DcuatroConclusion{}». Predomina el historial plano; los contrastes dentro del estudio son mayoritariamente positivos, con sesgo de notificación posible; el historial de evaluación a menudo no se reporta.
\textit{Limitaciones:} Autora única, revisores de IA de una misma familia, informes no recuperados, sesgo de notificación posible, ausencia de un conjunto de referencia independiente y solo inglés y español.
\textit{Conclusiones:} Con las salvedades del marco de elegibilidad, el escenario multiturno está menos investigado; se observan brechas de cobertura, no necesidades, en la evaluación con historial predicho y su reporte y en los datos en español. La combinación de estado estructurado con \en{joins} algorítmicos o retroalimentación de ejecución es poco frecuente, en línea con lo que cabría esperar por azar dadas las bajas prevalencias de cada elemento.
\textit{Financiamiento y registro:} Sin financiamiento específico; registro en OSF (DOI 10.17605/OSF.IO/5USF3), con enmiendas A4--A6 pendientes de actualización.
\end{abstract}

\begin{IEEEkeywords}
\en{Text-to-SQL} conversacional, revisión sistemática, vinculación de esquemas, estado del diálogo, autocorrección, PRISMA 2020
\end{IEEEkeywords}

\begin{abstracten}
\textit{Background:} In conversational Text-to-SQL each question depends on earlier turns, and no systematic review was identified that jointly analyzes context restoration, dialogue state, schema linking and self-correction.
\textit{Objective:} To classify multi-turn techniques, contrast them with single-turn ones, and synthesize their effect, evaluation and gaps.
\textit{Methods:} Systematic review following PRISMA 2020 with a registered protocol. The search covered \nBases{} sources (Scopus, IEEE Xplore, ACM Digital Library, ACL Anthology and arXiv) on \fechaBusqueda{}, plus snowballing. Empirical multi-turn studies were included, and single-turn studies only if schema linking or self-correction was their central contribution. AI models screened studies and extracted data under the author's verification; quality was assessed with a \nItemsCalidad{}-item Kitchenham-style checklist; synthesis was narrative (SWiM).
\textit{Results:} \nEstudios{} studies were included, \nMultiturno{} of them multi-turn. The multi-turn proportion (D4), descriptive and within the eligibility framework, was \DcuatroPct~\% (95~\% CI \DcuatroICinf{} to \DcuatroICsup~\%); the registered rule classifies it as ``considerably under-researched''. Flat history predominates; within-study contrasts are mostly positive, with possible reporting bias; the evaluation history is often not reported.
\textit{Limitations:} Single author, AI reviewers from one model family, unretrieved reports, possible reporting bias, no independent reference set, and English and Spanish only.
\textit{Conclusions:} Subject to the eligibility framework, the multi-turn setting is less researched; coverage gaps, not needs, are observed in evaluation with predicted history and its reporting, and in Spanish-language data. The combination of structured state with algorithmic joins or execution feedback is infrequent, in line with what chance would predict given the low prevalence of each element.
\textit{Funding and registration:} No specific funding; registered on OSF (DOI 10.17605/OSF.IO/5USF3), with amendments A4--A6 pending update.
\end{abstracten}

\begin{keywordsen}
Conversational Text-to-SQL, systematic review, schema linking, dialogue state, self-correction, PRISMA 2020
\end{keywordsen}
